% substitution-mstep.tex — Closed-form substitution M-steps for specific models
% To be \input from tkf.tex

\section{Substitution M-Steps for Specific Models}
\label{sec:subst-mstep}

This appendix specializes the general GTR substitution M-step
(Section~\ref{sec:bw-tkf91}) to specific DNA and codon substitution models.
In each case, we state the rate matrix $\revsub$ in terms of the model
parameters, derive the closed-form MLE (or MAP) update given the
bridge expectations $(W_i, U_{ij}, V_i)$,
and note the spectral structure when it simplifies the
endpoint-conditioned integrals~\eqref{eq:em-Jkl}.

Throughout, we work with the complete-data log-likelihood~\eqref{eq:ctmc-complete}
\[
\ell_2 = \sum_i V'_i \log \eqm_i
  + \sum_{j>i} (U_{ij} + U_{ji}) \log \exch_{ij}
  - \sum_{j>i} \exch_{ij}(W_i \eqm_j + W_j \eqm_i)
\]
where $V'_i = V_i + \sum_{j \neq i} U_{ji}$,
under the GTR parameterization $\revsub_{ij} = \exch_{ij} \eqm_j$
with symmetric exchangeability $\exch_{ij} = \exch_{ji}$.

The indel M-step (quadratic~\eqref{eq:kappa-quadratic} for $\kappa$,
with L'H\^opital limits in Section~\ref{sec:lhopital})
is orthogonal and unchanged across all substitution models.

For DNA models ($|\alphabet| = 4$), the bridge-expectation integral
$I^{ab}_{ij}(\evoltime)$~\eqref{eq:em-integral}
involves only $4 \times 4 = 16$ matrix entries per endpoint pair,
and the $J^{kl}(\evoltime)$ matrix~\eqref{eq:em-Jkl}
is $4 \times 4$,
making the endpoint-conditioned expectations cheap to compute.


%%%%%%%%%%%%%%%%%%%%%%%%%%%%%%%%%%%%%%%%%%%%%%%%%%%%%%%%%%%%%%%%%%%%%%%%
\subsection{JC69 (Jukes--Cantor)}
\label{sec:mstep-jc69}

\paragraph{Rate matrix.}
All off-diagonal rates are equal:
$\revsub_{ij} = \rho/3$ for $i \neq j$
(total rate out of any state is $\rho$).
Equivalently, $\eqm_i = 1/4$ for all $i$
and $\exch_{ij} = 4\rho/3$ for all $i \neq j$.

The single free parameter is the overall rate $\rho$.

\paragraph{M-step.}
By detailed balance,
$\exch_{ij} = (U_{ij} + U_{ji}) / (W_i \eqm_j + W_j \eqm_i)$
is the GTR MLE.
Under JC69, all exchangeabilities are forced equal, so we pool:
\begin{equation}
\label{eq:jc69-mstep}
\hat{\exch} = \frac{\sum_{j>i}(U_{ij} + U_{ji})}
       {\sum_{j>i}(W_i \eqm_j + W_j \eqm_i)}
= \frac{U_{\bullet}}{3 W_{\bullet}/4}
= \frac{4 U_{\bullet}}{3 W_{\bullet}}
\end{equation}
where $U_\bullet = \sum_{i \neq j} U_{ij}$ is the total substitution count
and $W_\bullet = \sum_i W_i$ is the total dwell time.
(The denominator uses $\sum_{j>i}(W_i/4 + W_j/4) = 3W_\bullet/4$
since each of the 4 states appears in 3 of the $\binom{4}{2} = 6$ pairs.)
The overall rate is then $\hat{\rho} = 3\hat{\exch}/4 = U_\bullet / W_\bullet$,
which is simply the total substitution count divided by total dwell time.

The equilibrium frequencies are fixed: $\hat{\eqm}_i = 1/4$.

\paragraph{Spectral structure.}
$\symsub = (\rho/3)J - (4\rho/3)\,I$ where $J$ is the $4 \times 4$ all-ones matrix.
Eigenvalues: $0$ (once, eigenvector $\mathbf{1}/2$) and $-4\rho/3$ (triply degenerate).
Thus $J^{kl}(\evoltime)$ has only two distinct diagonal values
($\evoltime$ and $\evoltime e^{-4\rho\evoltime/3}$)
and the off-diagonal $J^{01}$ term
$(1 - e^{-4\rho\evoltime/3})/(4\rho/3)$.


%%%%%%%%%%%%%%%%%%%%%%%%%%%%%%%%%%%%%%%%%%%%%%%%%%%%%%%%%%%%%%%%%%%%%%%%
\subsection{K80 (Kimura 2-Parameter)}
\label{sec:mstep-k80}

\paragraph{Rate matrix.}
Equal equilibrium frequencies $\eqm_i = 1/4$.
Transitions (purine$\leftrightarrow$purine, pyrimidine$\leftrightarrow$pyrimidine)
occur at rate $\kappa_s$ and transversions at rate $1$ (or vice versa,
depending on convention).
Labeling the states A, G, C, T and writing $i \sim j$ for a transition pair:
\[
\revsub_{ij} = \begin{cases}
  \kappa_s / 4 & \text{if } i \sim j \text{ (transition)} \\
  1/4 & \text{if } i \not\sim j \text{ (transversion)}
\end{cases}
\]
with total rate out of any state $(\kappa_s + 2)/4$.
The single free parameter beyond rate scaling is $\kappa_s$ (the
transition/transversion rate ratio).

\paragraph{M-step.}
Partition the substitution counts into transitions $U_s = \sum_{i \sim j} U_{ij}$
and transversions $U_v = \sum_{i \not\sim j} U_{ij} = U_\bullet - U_s$.
Similarly partition the dwell-weighted opportunities.
Under K80, the exchangeability for transition pairs is $\kappa_s$
and for transversion pairs is $1$ (up to a common scale).
The MLE for each group, from the pooled GTR formula:
\begin{align}
\hat{\kappa}_s &= \frac{U_s / N_s}{U_v / N_v}
\label{eq:k80-kappa}
\end{align}
where
$N_s = \sum_{i \sim j, \, i<j} (W_i \eqm_j + W_j \eqm_i)
     = W_\bullet / 4$
is the opportunity for transitions
(2 pairs: A--G and C--T, each contributing $(W_i+W_j)/4$,
summing to $W_\bullet/4$) and
$N_v = \sum_{i \not\sim j, \, i<j} (W_i \eqm_j + W_j \eqm_i)
     = W_\bullet / 2$
is the opportunity for transversions
(4 pairs, each state appearing in 2 transversion pairs,
summing to $2W_\bullet/4 = W_\bullet/2$).
Simplifying:
\begin{equation}
\label{eq:k80-mstep}
\hat{\kappa}_s = \frac{2\, U_s}{U_v}
\end{equation}
Given $\hat{\kappa}_s$, the overall scale $\hat\sigma$ satisfies
$\hat\sigma(\hat\kappa_s N_s + N_v) = U_\bullet$,
giving overall rate
$\hat\rho = \hat\sigma(\hat\kappa_s + 2)/4 = U_\bullet / W_\bullet$.

\paragraph{Spectral structure.}
Three distinct eigenvalues: $0$, $-(\kappa_s + 1)/2$ (doubly degenerate,
within-purine and within-pyrimidine contrasts), and $-1$
(purine-vs-pyrimidine contrast).


%%%%%%%%%%%%%%%%%%%%%%%%%%%%%%%%%%%%%%%%%%%%%%%%%%%%%%%%%%%%%%%%%%%%%%%%
\subsection{F81 (Felsenstein 1981)}
\label{sec:mstep-f81}

\paragraph{Rate matrix.}
Equal exchangeabilities, variable equilibrium:
$\revsub_{ij} = \rho\, \eqm_j$ for $i \neq j$
(equivalently, $\exch_{ij} = \rho$ for all $i \neq j$).

\paragraph{M-step.}
Since $\exch_{ij} = \rho$ is constant across all pairs,
the exchangeability MLE pools all pairs:
\begin{equation}
\label{eq:f81-rho}
\hat{\rho} = \frac{U_\bullet}{\sum_{j > i}(W_i \eqm_j + W_j \eqm_i)}
= \frac{U_\bullet}{\sum_i W_i (1 - \eqm_i)}
\end{equation}
using the identity
$\sum_{j>i}(W_i \eqm_j + W_j \eqm_i) = \sum_i W_i \sum_{j \neq i} \eqm_j = \sum_i W_i(1 - \eqm_i)$.

For equilibrium frequencies:
\begin{equation}
\label{eq:f81-pi}
\hat{\eqm}_i = \frac{V'_i}{\sum_j V'_j}
\end{equation}
This is the standard empirical-frequency estimator (exact MLE when
$\exch$ is held fixed; approximate when $\exch$ and $\eqm$ are jointly
optimized, as discussed in Section~\ref{sec:bw-tkf91}).

In practice, iterate: fix $\hat\eqm$, solve for $\hat\rho$~\eqref{eq:f81-rho},
update $\hat\eqm$~\eqref{eq:f81-pi}, repeat.
The decoupling is exact after one iteration when $W_i \propto \eqm_i$
(which holds approximately for long dwell times).

\paragraph{Spectral structure.}
$\symsub = \rho(\sqrt{\eqm}\sqrt{\eqm}^\top - I)$
where $\sqrt{\eqm}$ is the vector with entries $\sqrt{\eqm_i}$.
Eigenvalues: $0$ (once) and $-\rho$ ($|\alphabet|-1$ times, degenerate).
The $J^{kl}$ matrix has the same simple structure as JC69.


%%%%%%%%%%%%%%%%%%%%%%%%%%%%%%%%%%%%%%%%%%%%%%%%%%%%%%%%%%%%%%%%%%%%%%%%
\subsection{HKY85 (Hasegawa--Kishino--Yano)}
\label{sec:mstep-hky85}

\paragraph{Rate matrix.}
Variable equilibrium frequencies $\eqm$ and a
transition/transversion ratio $\kappa_s$:
\[
\revsub_{ij} = \begin{cases}
  \kappa_s\, \eqm_j & \text{if } i \sim j \text{ (transition)} \\
  \eqm_j & \text{if } i \not\sim j \text{ (transversion)}
\end{cases}
\]
Equivalently, $\exch_{ij} = \kappa_s$ for transition pairs and
$\exch_{ij} = 1$ for transversion pairs.
This is K80 + F81: the $\kappa_s$ structure of K80 with the
variable-$\eqm$ structure of F81.

\paragraph{M-step.}
Pool exchangeabilities within each class:
\begin{align}
\hat{\kappa}_s &= \frac{U_s \,/\, N_s}{U_v \,/\, N_v}
\label{eq:hky85-kappa}
\end{align}
where now (unlike K80) the opportunities depend on $\eqm$:
\begin{align}
N_s &= \sum_{i \sim j,\, i < j} (W_i \eqm_j + W_j \eqm_i) \label{eq:hky85-Ns} \\
N_v &= \sum_{i \not\sim j,\, i < j} (W_i \eqm_j + W_j \eqm_i) \label{eq:hky85-Nv}
\end{align}
With $U_s$ and $U_v$ as in K80 (summed over directed transition and
transversion pairs respectively).

For $\eqm$, use $\hat{\eqm}_i = V'_i / \sum_j V'_j$ as in F81,
then iterate with $\kappa_s$ if desired.

\paragraph{Spectral structure.}
Let $\eqm_R = \eqm_A + \eqm_G$ and $\eqm_Y = \eqm_C + \eqm_T$.
The eigenvalues are
$0$,
$-(\eqm_R + \kappa_s \eqm_Y)$ (within-pyrimidine contrast),
$-(\eqm_Y + \kappa_s \eqm_R)$ (within-purine contrast),
and $-1$ (purine-vs-pyrimidine contrast).
Generically four distinct values;
when $\eqm_R = \eqm_Y = 1/2$ (i.e.\ K80),
the two within-class eigenvalues merge to $-(\kappa_s + 1)/2$.


%%%%%%%%%%%%%%%%%%%%%%%%%%%%%%%%%%%%%%%%%%%%%%%%%%%%%%%%%%%%%%%%%%%%%%%%
\subsection{GTR (General Time-Reversible)}
\label{sec:mstep-gtr}

For completeness, we restate the general case from
Section~\ref{sec:bw-tkf91}.

\paragraph{Rate matrix.}
$\revsub_{ij} = \exch_{ij}\,\eqm_j$ with $\exch_{ij} = \exch_{ji} \geq 0$
for $i \neq j$.
For DNA, $\exch$ has 6 free parameters (upper triangle)
and $\eqm$ has 3 free parameters (4 frequencies summing to 1).

\paragraph{M-step.}
Fixing $\eqm$:
\begin{equation}
\label{eq:gtr-Q}
\hat{\exch}_{ij} = \frac{U_{ij} + U_{ji}}{W_i \eqm_j + W_j \eqm_i}
\end{equation}
Fixing $\exch$, the $\eqm$ update requires solving a nonlinear system
(Section~\ref{sec:bw-tkf91}).
In practice, set $\hat{\eqm}_i \propto V'_i$
and iterate with~\eqref{eq:gtr-Q}.

\paragraph{Spectral structure.}
$\symsub_{ij} = \exch_{ij}\sqrt{\eqm_i \eqm_j}$ is real symmetric
with eigenvalues $0 = \xi^{(0)} > \xi^{(1)} \geq \cdots \geq \xi^{(|\alphabet|-1)}$
(generically all distinct for DNA).
No further simplification beyond the general formulas~\eqref{eq:em-Jkl}.


%%%%%%%%%%%%%%%%%%%%%%%%%%%%%%%%%%%%%%%%%%%%%%%%%%%%%%%%%%%%%%%%%%%%%%%%
\subsection{GY94 (Goldman--Yang Codon Model)}
\label{sec:mstep-gy94}

\paragraph{Rate matrix.}
The state space is the 61 sense codons ($|\alphabet| = 61$,
excluding stop codons).
The instantaneous rate from codon $i$ to codon $j$ is:
\[
\revsub_{ij} = \begin{cases}
  0 & \text{if $i$ and $j$ differ at more than one position} \\
  \eqm_j & \text{if synonymous transversion} \\
  \kappa_s\, \eqm_j & \text{if synonymous transition} \\
  \omega\, \eqm_j & \text{if nonsynonymous transversion} \\
  \omega\, \kappa_s\, \eqm_j & \text{if nonsynonymous transition}
\end{cases}
\]
where $\omega$ is the nonsynonymous/synonymous rate ratio (dN/dS),
$\kappa_s$ is the transition/transversion ratio,
and $\eqm_j$ is the equilibrium frequency of codon $j$.
This is a reversible model with exchangeability
\[
\exch_{ij} = \begin{cases}
  0 & \text{(multi-nucleotide change)} \\
  1 & \text{(synonymous transversion)} \\
  \kappa_s & \text{(synonymous transition)} \\
  \omega & \text{(nonsynonymous transversion)} \\
  \omega\kappa_s & \text{(nonsynonymous transition)}
\end{cases}
\]

\paragraph{M-step.}
Partition the substitution counts into four classes indexed by
(synonymous/nonsynonymous) $\times$ (transition/transversion):
\begin{align*}
U_{S,s} &= \text{total synonymous transition counts} &
U_{S,v} &= \text{total synonymous transversion counts} \\
U_{N,s} &= \text{total nonsynonymous transition counts} &
U_{N,v} &= \text{total nonsynonymous transversion counts}
\end{align*}
and the corresponding dwell-weighted opportunities $N_{S,s}, N_{S,v}, N_{N,s}, N_{N,v}$
defined analogously to~\eqref{eq:hky85-Ns}--\eqref{eq:hky85-Nv}
but summing over single-nucleotide codon pairs in each class.

The exchangeabilities are $\exch = 1$ (synonymous transversion),
$\kappa_s$ (synonymous transition),
$\omega$ (nonsynonymous transversion),
$\omega\kappa_s$ (nonsynonymous transition).
The complete-data log-likelihood for $(\omega, \kappa_s)$,
with $\eqm$ fixed, is:
\begin{align*}
\ell_2(\omega,\kappa_s) &=
  U_{S,v}\log 1 + U_{S,s}\log\kappa_s + U_{N,v}\log\omega + U_{N,s}\log(\omega\kappa_s) \\
  &\quad - N_{S,v} - \kappa_s N_{S,s} - \omega N_{N,v} - \omega\kappa_s N_{N,s}
\end{align*}
Setting $\partial\ell_2/\partial\kappa_s = 0$ and $\partial\ell_2/\partial\omega = 0$:
\begin{align}
\frac{U_{S,s} + U_{N,s}}{\kappa_s} &= N_{S,s} + \omega\, N_{N,s} \label{eq:gy94-kappa-eq} \\
\frac{U_{N,v} + U_{N,s}}{\omega} &= N_{N,v} + \kappa_s\, N_{N,s} \label{eq:gy94-omega-eq}
\end{align}
These are two equations in two unknowns. Substituting~\eqref{eq:gy94-omega-eq}
into~\eqref{eq:gy94-kappa-eq} to eliminate $\omega$:
\[
\omega = \frac{U_{N,v} + U_{N,s}}{N_{N,v} + \kappa_s N_{N,s}}
\]
\[
\frac{U_{S,s} + U_{N,s}}{\kappa_s} = N_{S,s} + \frac{(U_{N,v} + U_{N,s})\, N_{N,s}}{N_{N,v} + \kappa_s N_{N,s}}
\]
Multiplying both sides by $\kappa_s(N_{N,v} + \kappa_s N_{N,s})$ and rearranging:
\begin{equation}
\label{eq:gy94-kappa-quad}
N_{S,s}\, N_{N,s}\;\kappa_s^2
+ \bigl(N_{S,s}\, N_{N,v} + (U_{N,v} - U_{S,s})\,N_{N,s}\bigr)\;\kappa_s
- (U_{S,s}+U_{N,s})\, N_{N,v} = 0
\end{equation}
The positive root gives $\hat{\kappa}_s$,
and then
\begin{equation}
\label{eq:gy94-omega}
\hat{\omega} = \frac{U_{N,v} + U_{N,s}}{N_{N,v} + \hat{\kappa}_s\, N_{N,s}}
\end{equation}

The codon equilibrium frequencies $\eqm$ are estimated as
$\hat{\eqm}_j \propto V'_j$.

\paragraph{Spectral structure.}
For 61 sense codons, the symmetrized rate matrix is $61 \times 61$
with at most 61 distinct eigenvalues.
No closed-form eigendecomposition is known in general for
arbitrary $\eqm$.
When $\eqm$ factors as $\eqm_{abc} = \eqm^{(1)}_a \eqm^{(2)}_b \eqm^{(3)}_c$
(the F3x4 model), the symmetrized matrix is a sum of Kronecker products
and its eigenvectors factor as tensor products of $4\times 4$ matrices,
reducing the eigendecomposition to three $4\times 4$ problems.
In practice, for the full GY94, numerical eigendecomposition of the
$61 \times 61$ symmetric matrix is inexpensive ($O(61^3) \approx 2 \times 10^5$ flops)
and need only be performed once per M-step.


%%%%%%%%%%%%%%%%%%%%%%%%%%%%%%%%%%%%%%%%%%%%%%%%%%%%%%%%%%%%%%%%%%%%%%%%
\subsection{Summary and Practical Notes}
\label{sec:subst-mstep-summary}

\begin{center}
\begin{tabular}{lccl}
\toprule
Model & Free params & $\eqm$ & M-step for exchangeability \\
\midrule
JC69 & 1 ($\rho$) & fixed $1/4$ & $\hat\rho = U_\bullet/W_\bullet$ \\
K80 & 1 ($\kappa_s$) & fixed $1/4$ & $\hat\kappa_s = 2U_s/U_v$ \\
F81 & 4 ($\rho, \eqm$) & $V'_i/\sum V'_j$ & $\hat\rho = U_\bullet / \sum_i W_i(1-\eqm_i)$ \\
HKY85 & 5 ($\kappa_s, \eqm$) & $V'_i/\sum V'_j$ & $\hat\kappa_s = (U_s/N_s)/(U_v/N_v)$ \\
GTR & 9 ($\exch, \eqm$) & $V'_i/\sum V'_j$ & $\hat\exch_{ij} = (U_{ij}+U_{ji})/(W_i\eqm_j+W_j\eqm_i)$ \\
GY94 & 62 ($\omega,\kappa_s,\eqm$) & $V'_j/\sum V'_k$ & quadratic in $\kappa_s$~\eqref{eq:gy94-kappa-quad} \\
\bottomrule
\end{tabular}
\end{center}

For all models above, the indel M-step is identical: solve
the $\kappa$ quadratic~\eqref{eq:kappa-quadratic} for
$(\insrate, \delrate)$.
The substitution and indel parameter blocks decouple
completely in the M-step (they share no sufficient statistics).

\paragraph{MAP estimates.}
To obtain MAP estimates,
augment the sufficient statistics before applying the MLE formulas:
$U_{ij} \to U_{ij} + \alpha_\exch - 1$,
$W_i \to W_i + \beta_\exch$,
$V_i \to V_i + \alpha_\eqm - 1$\ifdefined\MIXDOMPAPER ,
as in Section~\ref{sec:bw-mixdom}\fi.
The independent Gamma$(\alpha_\exch,\beta_\exch)$ priors on each
$\exch_{ij}$ together with a Dirichlet on $\eqm$ are conjugate to the
\emph{irreversible} CTMC complete-data likelihood (independent
$\revsub_{ij}$) but not to the reversible GTR
parameterization---the symmetric coupling $\exch_{ij} = \exch_{ji}$
together with the requirement that the joint complete-data likelihood
include the stationary draw $X(0) \sim \eqm$ means the proper
conjugate prior is the cycle-corrected edge-flow density of
Diaconis and Rolles~\cite{DiaconisRolles2006}, with a Kirchhoff
matrix-tree determinant accounting for the cycle structure of the
state graph.\ifdefined\MIXDOMPAPER{}  See the discussion in Section~\ref{sec:bw-mixdom-priors}
for the form of the Diaconis--Rolles prior and its modification to
incorporate the initial-state contribution.\fi  The simpler
Gamma${}\times{}$Dirichlet pseudocounts used here are non-conjugate
regularizers and remain perfectly valid for MAP, with the closed-form
augmented M-step above; we adopt them throughout for simplicity.

\paragraph{Degenerate eigenvalues.}
Models with symmetry (JC69, K80, F81) have degenerate eigenvalues
in $\symsub$, which means the $J^{kl}(\evoltime)$ matrix has
repeated diagonal entries.
This reduces the number of distinct $J$ values that must be computed,
but requires care to use the $\xi^{(k)} = \xi^{(l)}$ branch
($J^{kl} = \evoltime\, e^{\xi^{(k)}\evoltime}$)
rather than the difference quotient,
which is a removable $0/0$ singularity.


%%%%%%%%%%%%%%%%%%%%%%%%%%%%%%%%%%%%%%%%%%%%%%%%%%%%%%%%%%%%%%%%%%%%%%%%
\subsection{Reversible Mixture with Per-Component GTR}
\label{sec:mstep-mixture}

We now consider a $C$-component reversible mixture in which each
component is a full GTR model with its own exchangeability matrix
$\exch^{(c)}$ and equilibrium distribution $\eqm^{(c)}$:
\begin{equation}\label{eq:mixture-Q}
\revsub^{(c)}_{ij} = \exch^{(c)}_{ij}\,\eqm^{(c)}_j
\quad (i \neq j),
\qquad c \in \{1,\dots,C\}.
\end{equation}
Each site is assigned (latently) to one component, and conditional
on $c$, the site evolves under the reversible CTMC with rate
matrix~$\revsub^{(c)}$.
\ifdefined\MIXDOMPAPER This is the parameterisation used by the per-(domain, fragment-type)
site-class emissions of MixDom (Section~\ref{sec:bw-mixdom}), where
$c$ is the site-class drawn independently at each position from
$\classdist_{\dom\frag}$.\else This is the parameterisation used by per-site-class emissions in
mixture-of-CTMC models, where $c$ is the latent site class drawn
independently at each position from a per-site Dirichlet.\fi
Because each component is an independent GTR model, the overall
rate scale is absorbed into $\exch^{(c)}$ itself and no separate
rate multiplier is needed; this makes the parameterisation a
strict generalisation of per-domain GTR (the per-domain GTR case
is recovered when $\classdist_{\dom\frag,c} = \mathbf{1}[c = \dom]$)
and of the
discrete-gamma rate-across-sites model~\cite{Yang1994a}.

\paragraph{E-step counts.}
The component assignment is latent. From the standard E-step,
each branch produces an endpoint pair $(a,b)$ with marginal posterior
weight $w_{ab}$; conditional on $(a,b)$ at a site with prior class
distribution $P(c \mid \text{site}) = \classdist_c$, the
component posterior is
\begin{equation}\label{eq:mix-post}
P(c \mid a,b,\evoltime,\text{site})
  = \frac{\classdist_c\,\eqm^{(c)}_a\,M^{(c)}_{ab}(\evoltime)}
         {\sum_{c'} \classdist_{c'}\,\eqm^{(c')}_a\,M^{(c')}_{ab}(\evoltime)},
\end{equation}
with $M^{(c)}(\evoltime) = \exp(\revsub^{(c)}\evoltime)$.
Per-component bridge-expectation sufficient statistics are then
accumulated by weighting the standard
expressions~\eqref{eq:em-counts}--\eqref{eq:em-dwell} by~\eqref{eq:mix-post}:
\begin{align}
V^{(c)}_i &= \textstyle\sum_{\text{starts}} w_a\,
   P(c\mid a,\cdot)\,\mathbf{1}[X(0)=i] \label{eq:mix-V}\\
U^{(c)}_{ij} &= \textstyle\sum_{\text{branches}} w_{ab}\,
   P(c\mid a,b,\evoltime,\cdot)\,
   \emcount^U_{ij}\!\bigl(a,b,\evoltime;\,\revsub^{(c)}\bigr) \label{eq:mix-U}\\
W^{(c)}_i &= \textstyle\sum_{\text{branches}} w_{ab}\,
   P(c\mid a,b,\evoltime,\cdot)\,
   \emcount^W_i\!\bigl(a,b,\evoltime;\,\revsub^{(c)}\bigr) \label{eq:mix-W}
\end{align}
and we write $V'^{(c)}_i = V^{(c)}_i + \sum_{j \neq i} U^{(c)}_{ji}$ as before.
Writing $U^{(c)}_\bullet = \sum_{i\neq j} U^{(c)}_{ij}$ for the
total per-component substitution count.

\paragraph{Complete-data log-likelihood.}
Because each class is an independent GTR model under this
parameterisation, the complete-data log-likelihood
decomposes over classes:
\begin{equation}\label{eq:mix-loglik}
\ell_2(\{\exch^{(c)}\},\{\eqm^{(c)}\}) = \sum_c \ell_2^{(c)},
\end{equation}
\begin{align}
\ell_2^{(c)}
&= \sum_i V'^{(c)}_i \log \eqm^{(c)}_i
+ \sum_{j>i}\bigl(U^{(c)}_{ij}+U^{(c)}_{ji}\bigr)\log \exch^{(c)}_{ij}
\notag\\
&\quad - \sum_{j>i} \exch^{(c)}_{ij}
   \bigl(W^{(c)}_i \eqm^{(c)}_j + W^{(c)}_j \eqm^{(c)}_i\bigr).
\label{eq:mix-loglik-c}
\end{align}

\paragraph{Identifiability.}
No gauge to fix: the rate scale of class $c$ lives inside
$\exch^{(c)}$, and the $\{\exch^{(c)}\}$ are uncoupled in both
\eqref{eq:mixture-Q} and~\eqref{eq:mix-loglik-c}.
Every class can be updated independently.

\paragraph{Per-class GTR M-step.}
Since $\ell_2$ decomposes class by class, the update for class $c$ is
exactly the GTR M-step of Section~\ref{sec:bw-tkf91} applied to that
class's sufficient statistics.
Differentiating~\eqref{eq:mix-loglik-c} w.r.t.\ $\exch^{(c)}_{ij}$:
\[
\frac{\partial \ell_2^{(c)}}{\partial \exch^{(c)}_{ij}}
= \frac{U^{(c)}_{ij}+U^{(c)}_{ji}}{\exch^{(c)}_{ij}}
- \bigl(W^{(c)}_i \eqm^{(c)}_j + W^{(c)}_j \eqm^{(c)}_i\bigr) = 0,
\]
giving the closed-form update
\begin{equation}\label{eq:mix-S-update}
\hat\exch^{(c)}_{ij}
= \frac{U^{(c)}_{ij} + U^{(c)}_{ji}}
       {W^{(c)}_i \eqm^{(c)}_j + W^{(c)}_j \eqm^{(c)}_i}.
\end{equation}
The equilibrium update uses the empirical-frequency approximation
\begin{equation}\label{eq:mix-pi-update}
\hat\eqm^{(c)}_i \propto V'^{(c)}_i,
\end{equation}
which is exact when $\exch^{(c)}_{ij} W^{(c)}_j$ is independent of
$i$ (F81 limit) and approximate otherwise, identical to
Section~\ref{sec:bw-tkf91}.
No alternation or projection is needed; one pass through each class
completes the M-step.

\paragraph{Connection to standard models.}
With $C=1$, \eqref{eq:mix-S-update}--\eqref{eq:mix-pi-update} reduce
exactly to the GTR M-step~\eqref{eq:gtr-Q}.
With $\classdist_{\dom\frag,c} = \mathbf{1}[c=\dom]$ so each domain
uses its own dedicated class, the model is the per-domain GTR
($\exch^{(\dom)}$) special case; the per-class M-step here reduces to
the per-domain M-step of Section~\ref{sec:bw-tkf91} applied once per
domain.
With $\eqm^{(c)} \equiv \eqm$ and $\exch^{(c)} = \gamma_c\,\exch$ for
some shared $\exch$, the model is the discrete-gamma rate
model~\cite{Yang1994a}; a constrained M-step that enforces
$\exch^{(c)} \propto \exch$ recovers the shared-exchangeability case
as a special case.
Free-rate variants~\cite{Susko2018} are the shared-$\exch$ /
free-$\gamma_c$ special case.

\paragraph{MAP estimates.}
Per-class regularising priors apply independently to each class.
A Dirichlet pseudocount $\alpha^{(c)}_i$ is added to each
$V^{(c)}_i$, and a Gamma$(a_S, b_S)$ prior on each
$\exch^{(c)}_{ij}$ (e.g.\ calibrated to LG via
$a_S = b_S\,\exch^{\mathrm{LG}}_{ij}$) adds $a_S-1$ to the numerator
and $b_S$ to the denominator of~\eqref{eq:mix-S-update}.
As above, this independent-Gamma${}\times{}$Dirichlet structure is
conjugate only in the irreversible parameterization; it is a
non-conjugate regularizer for the reversible GTR mixture component,
with the proper conjugate being a per-class Diaconis--Rolles
prior~\cite{DiaconisRolles2006}.  We use the simpler form here for
the same reasons as the single-component case\ifdefined\MIXDOMPAPER{}
(Section~\ref{sec:bw-mixdom-priors})\fi.
The MAP objective
\begin{equation}\label{eq:mix-map-obj}
\mathcal{L} = \ell_2
  + \sum_{c,i}(\alpha^{(c)}_i - 1)\log\eqm^{(c)}_i
  + \sum_{c,j>i}\bigl[(a_S - 1)\log\exch^{(c)}_{ij}
          - b_S \exch^{(c)}_{ij}\bigr]
\end{equation}
decomposes across classes and has a \emph{unique} joint maximum
at
\begin{equation}\label{eq:mix-map-S}
\hat\exch^{(c)}_{ij}
= \frac{U^{(c)}_{ij}+U^{(c)}_{ji} + a_S - 1}
       {W^{(c)}_i \eqm^{(c)}_j + W^{(c)}_j \eqm^{(c)}_i + b_S},
\end{equation}
with the $\eqm^{(c)}$ update obtained by the same bisection as
Section~\ref{sec:bw-tkf91}.
Each class is a standalone GTR M-step; no cross-class coupling.

\paragraph{Hyperparameter sizing under SVI.}
Under SVI-BW~(\ref{eq:svi-update}) the EMA sufficient statistics are
scaled to the full-dataset size~$N$, so prior pseudocounts must
live on the same scale to retain their relative weight.
Parameterising by \emph{effective sample size},
set $\alpha^{(c)}_i = N_\pi\,\eqm^{\mathrm{LG}}_i$
and $b_S = N_S$ with $a_S = N_S\,\exch^{\mathrm{LG}}_{ij}$.
Defaults $N_\pi, N_S \sim 10^2\text{--}10^3$
(data-independent) give priors that vanish in the data limit
but identify the posterior in small batches.


%%%%%%%%%%%%%%%%%%%%%%%%%%%%%%%%%%%%%%%%%%%%%%%%%%%%%%%%%%%%%%%%%%%%%%%%
\subsection{Rate Rescaling: M-Step for a Global Scalar Multiplier}
\label{sec:mstep-rescaling}

In some training scenarios the relative pattern of substitution rates
is known (e.g.\ from a previously-estimated GTR matrix or from an
external published matrix such as LG) and only the overall \emph{rate}
needs to adapt to the current data.  Concretely, we hold both the
exchangeability matrix~$\exch$ \emph{and} the equilibrium
distribution~$\eqm$ fixed, and seek a single scalar multiplier
$\sigma > 0$ such that the rescaled rate matrix
\begin{equation}\label{eq:rescale-Q}
\revsub^{\mathrm{new}}_{ij} = \sigma\,\exch_{ij}\,\eqm_j
\quad (i \neq j)
\end{equation}
maximises the complete-data log-likelihood~$\ell_2$.
Equivalently, the new exchangeability is
$\exch^{\mathrm{new}}_{ij} = \sigma\,\exch_{ij}$,
so all relative exchangeabilities are preserved and only the overall
tempo of the model changes.

\paragraph{Closed-form $\sigma$.}
Substituting $\exch_{ij} \to \sigma\exch_{ij}$ into the
GTR complete-data log-likelihood,
\[
\ell_2(\sigma) = \mathrm{const}
  + \log\sigma\;\sum_{j>i}(U_{ij} + U_{ji})
  - \sigma \sum_{j>i}\exch_{ij}\bigl(W_i\eqm_j + W_j\eqm_i\bigr).
\]
The terms not involving $\sigma$ (the $V'_i\log\eqm_i$ part and
$\sum_{j>i}(U_{ij}+U_{ji})\log\exch_{ij}$) are constants under
rescaling.  Differentiating and setting to zero yields the unique
maximiser
\begin{equation}\label{eq:rescale-sigma}
\hat\sigma = \frac{U_\bullet}{D},
\qquad
U_\bullet = \sum_{i\neq j} U_{ij},
\qquad
D = \sum_{j>i}\exch_{ij}\bigl(W_i\eqm_j + W_j\eqm_i\bigr),
\end{equation}
where $U_\bullet$ is the total expected substitution count and $D$ is
the total dwell-weighted opportunity for substitution under the fixed
shape $(\exch, \eqm)$.  The objective is strictly concave in $\sigma$
on $(0, \infty)$ (the Hessian $-U_\bullet/\sigma^2 < 0$), so
\eqref{eq:rescale-sigma} is the global maximiser.

\paragraph{Mixture form (per-class).}
Applied class by class to the reversible mixture of
Section~\ref{sec:mstep-mixture}, each class $c$ has its own scalar
$\sigma_c$ updated independently from its own per-class statistics:
\begin{equation}\label{eq:rescale-sigma-c}
\hat\sigma_c
= \frac{U^{(c)}_\bullet}{D^{(c)}},
\qquad
D^{(c)} = \sum_{j>i}\exch^{(c)}_{ij}
   \bigl(W^{(c)}_i \eqm^{(c)}_j + W^{(c)}_j \eqm^{(c)}_i\bigr).
\end{equation}
The shapes $\exch^{(c)}$ and $\eqm^{(c)}$ are held at their current
values and only the per-class rate scale $\sigma_c$ moves; the
class-decomposition~\eqref{eq:mix-loglik} ensures classes do not couple.
A common-rate variant (a single shared $\hat\sigma$ across all classes)
is recovered by pooling: $\hat\sigma = \sum_c U^{(c)}_\bullet
/ \sum_c D^{(c)}$.  The per-class form~\eqref{eq:rescale-sigma-c}
attains a strictly higher $\ell_2$ unless all $\sigma_c$ already coincide.

\paragraph{Aggregation of counts.}
The only sufficient statistics required are
$U^{(c)}_\bullet = \sum_{i\neq j} U^{(c)}_{ij}$ (a scalar per class)
and the dwell--frequency contraction~$D^{(c)}$, which is itself a
scalar inner product of the fixed quantity
$\exch^{(c)}_{ij}(W^{(c)}_i \eqm^{(c)}_j + W^{(c)}_j \eqm^{(c)}_i)$
against the per-pair count.  Both are linear in the underlying
per-pair $(W, U)$ statistics, so under SVI-BW they accumulate via the
same EMA used for the standard M-step (no special handling).

\paragraph{MAP estimate.}
With a Gamma$(a_\sigma, b_\sigma)$ prior on $\sigma$, the augmented
M-step is
\begin{equation}\label{eq:rescale-map}
\hat\sigma = \frac{U_\bullet + a_\sigma - 1}{D + b_\sigma}.
\end{equation}
The Gamma prior on the global rate is conjugate (the model is linear
in $\sigma$), in contrast to the non-conjugate prior on
the full-shape M-step.  In the rate-rescaling regime the standard
$(N_\pi, N_S)$ pseudocounts on $(\exch, \eqm)$ are inactive (those
parameters are frozen); $(a_\sigma, b_\sigma)$ centred at $\sigma = 1$
(e.g.\ $a_\sigma = b_\sigma = N_\sigma$ for some effective sample size
$N_\sigma$) regularises the rate towards the current frozen scale.

\paragraph{When to use.}
Rate rescaling is a useful intermediate regime between fully-frozen
substitution parameters (no M-step) and the full-shape mixture M-step
of Section~\ref{sec:mstep-mixture}.  It is particularly natural when
warm-starting from a published matrix and one wishes to absorb a
calibration drift in the time units of the data \emph{without}
re-estimating the relative chemistry.  In an alternating schedule
(e.g.\ alternate $\sigma_c$ updates with tied-$\eqm$
\eqref{eq:tied-pi-pi-update} updates) the rescaling step
contributes a single guaranteed $\ell_2$ ascent per class per
iteration, and the alternation does not destabilise the joint
objective because each step is an exact maximisation in its own
coordinate.


%%%%%%%%%%%%%%%%%%%%%%%%%%%%%%%%%%%%%%%%%%%%%%%%%%%%%%%%%%%%%%%%%%%%%%%%
\subsection{Joint Rate Rescaling and Equilibrium}
\label{sec:mstep-rescaling-pi}

A natural extension of the rate-rescaling regime
(Section~\ref{sec:mstep-rescaling}) frees the equilibrium
distribution~$\eqm$ jointly with the scale~$\sigma$ while keeping the
exchangeability shape $\exch$ fixed.  This matches the substitution
degrees of freedom of a Maraschino-style fit with
\verb|--freeze-class-S-shape| (per-class $\sigma_c$ and $\eqm^{(c)}$
free, $\exch^{(c)}$ shape locked at LG), and is the natural
SVI--BW counterpart when warm-starting from such a fit: the previous
\verb|rescaling-rates| mode artificially froze $\eqm^{(c)}$ at the
warm-start values, dropping $A-1$ degrees of freedom per class.

\paragraph{Joint complete-data log-likelihood.}
With $\revsub_{ij} = \sigma\,\exch_{ij}\,\eqm_j$ and $\exch$ frozen,
the complete-data log-likelihood as a function of $(\sigma, \eqm)$ is
\begin{equation}\label{eq:rescale-pi-ll}
\ell_2(\sigma, \eqm) =
   \sum_i V'_i \log \eqm_i + U_\bullet \log \sigma
   - \sigma \sum_{j>i} \exch_{ij}\bigl(W_i \eqm_j + W_j \eqm_i\bigr)
   + \mathrm{const},
\end{equation}
where the constant collects the frozen $\sum_{j>i}(U_{ij}+U_{ji})\log\exch_{ij}$
term.  Using $\exch_{ij} = \exch_{ji}$ and $\exch_{ii} = 0$, the dwell
penalty contracts as
\[
\sum_{j>i}\exch_{ij}\bigl(W_i\eqm_j + W_j\eqm_i\bigr)
= \sum_b \eqm_b\,r_b,
\qquad
r_b := \sum_{a\neq b} \exch_{ab} W_a = (\exch W)_b,
\]
so $r_b$ depends only on the frozen quantities $(\exch, W)$ and is
constant under variation of $(\sigma, \eqm)$.  Equation
\eqref{eq:rescale-pi-ll} simplifies to
\begin{equation}\label{eq:rescale-pi-ll-simple}
\ell_2(\sigma, \eqm) =
   \sum_b V'_b \log \eqm_b + U_\bullet \log \sigma
   - \sigma\,\mathbf{r}^\top \eqm + \mathrm{const}.
\end{equation}

\paragraph{Stationary conditions.}
Lagrangian for $\sum_b \eqm_b = 1$ with multiplier $\lambda$:
\[
\mathcal{L}(\sigma,\eqm,\lambda)
= \ell_2(\sigma,\eqm) - \lambda\Bigl(\sum_b \eqm_b - 1\Bigr).
\]
The first-order conditions are
\begin{align}
\frac{\partial \mathcal{L}}{\partial \sigma}
   &= \frac{U_\bullet}{\sigma} - \mathbf{r}^\top\eqm = 0
   \;\Longrightarrow\;
   \sigma = \frac{U_\bullet}{\mathbf{r}^\top\eqm},
   \label{eq:rescale-pi-foc-sigma}\\[2pt]
\frac{\partial \mathcal{L}}{\partial \eqm_b}
   &= \frac{V'_b}{\eqm_b} - \sigma r_b - \lambda = 0
   \;\Longrightarrow\;
   \eqm_b = \frac{V'_b}{\lambda + \sigma r_b}.
   \label{eq:rescale-pi-foc-pi}
\end{align}

\paragraph{The Lagrange multiplier is fixed.}
A clean structural property obtains by multiplying
\eqref{eq:rescale-pi-foc-pi} through by $\eqm_b$ and summing:
\[
\sum_b V'_b = \lambda \sum_b \eqm_b + \sigma\, \mathbf{r}^\top \eqm
= \lambda + \sigma\,\mathbf{r}^\top\eqm.
\]
The $\sigma$-FOC \eqref{eq:rescale-pi-foc-sigma} sets
$\sigma\,\mathbf{r}^\top\eqm = U_\bullet$, so
\[
N \;=\; \lambda + U_\bullet,
\qquad N := \sum_b V'_b = V_\bullet + U_\bullet,
\]
which gives
\begin{equation}\label{eq:rescale-pi-lambda}
\boxed{\;\lambda \;=\; V_\bullet,\;}
\end{equation}
\emph{independent of $\sigma$}: the multiplier is fully determined by
the boundary count $V_\bullet = \sum_i V_i$.

\paragraph{Closed-form $\eqm(\sigma)$ and 1-D root in $\sigma$.}
Substituting~\eqref{eq:rescale-pi-lambda}
into~\eqref{eq:rescale-pi-foc-pi},
\begin{equation}\label{eq:rescale-pi-pi-of-sigma}
\hat\eqm_b(\sigma) = \frac{V'_b}{V_\bullet + \sigma\, r_b}.
\end{equation}
Note $\sum_b \hat\eqm_b(\sigma) = 1$ at $\sigma$ satisfying the
$\sigma$-FOC, automatically.  Plugging back into the
$\sigma$-FOC \eqref{eq:rescale-pi-foc-sigma}, $\sigma$ satisfies the
single one-dimensional equation
\begin{equation}\label{eq:rescale-pi-sigma-eq}
g(\sigma) := \sum_b \frac{\sigma\, r_b\, V'_b}{V_\bullet + \sigma\, r_b}
\;-\; U_\bullet \;=\; 0,
\qquad \sigma > 0.
\end{equation}

\paragraph{Existence and uniqueness.}
$g$ is well-defined and smooth for $\sigma > 0$ (denominators are
positive when $V_\bullet > 0$ or $r_b > 0$).  Boundary behaviour:
\[
g(0) = -U_\bullet \le 0,
\qquad
g(\sigma) \xrightarrow[\sigma\to\infty]{} \sum_b V'_b - U_\bullet
= V_\bullet \ge 0,
\]
and the derivative is strictly positive,
\[
g'(\sigma) = \sum_b \frac{r_b V'_b\, V_\bullet}{(V_\bullet + \sigma r_b)^2} > 0
\quad\text{whenever some } r_b V'_b > 0.
\]
Hence $g$ is strictly increasing and crosses zero exactly once at a
unique $\hat\sigma > 0$ (provided $V_\bullet > 0$ and $U_\bullet > 0$;
otherwise the boundary $\hat\sigma = 0$ or $\hat\sigma \to \infty$ is
the maximiser and the regime is degenerate).  The joint stationary
point $(\hat\sigma, \hat\eqm)$ is thus globally unique on the simplex.

\paragraph{Numerical solution: 1-D Newton.}
$g$ is monotone with a closed-form derivative, so a few Newton steps
on $\log\sigma$ (or bisection in $\sigma$) starting from the warm
$\sigma_{\mathrm{old}}$ converges to machine precision in $\le 10$
iterations.  Once $\hat\sigma$ is found, $\hat\eqm$ follows from
\eqref{eq:rescale-pi-pi-of-sigma}.

\paragraph{Equivalent: coordinate ascent.}
Both individual updates~\eqref{eq:rescale-pi-foc-sigma} and
\eqref{eq:rescale-pi-foc-pi} are exact closed-form maximisations of
$\ell_2$ in their own coordinate (for fixed value of the other),
so alternating
\begin{align*}
\sigma &\leftarrow U_\bullet \,/\, (\mathbf{r}^\top \eqm),\\
\eqm_b &\leftarrow V'_b \,/\, (V_\bullet + \sigma\, r_b)
\quad (\text{automatically normalised once $\lambda = V_\bullet$ holds})
\end{align*}
gives a monotone $\ell_2$ ascent and converges to the unique fixed
point identified above.  In code, $\sim 5$ alternation rounds
typically suffice; the 1-D Newton on \eqref{eq:rescale-pi-sigma-eq} is
slightly faster and gives the same answer.

\paragraph{Mixture form (per-class).}
Applied class by class to the reversible mixture of
Section~\ref{sec:mstep-mixture}, each class $c$ has its own pair
$(\sigma_c, \eqm^{(c)})$ updated independently from its own per-class
sufficient statistics $(W^{(c)}, U^{(c)}, V^{(c)})$ and its own fixed
$\exch^{(c)}$ shape:
\begin{equation}\label{eq:rescale-pi-class}
\hat\sigma_c\;\text{solves}\;
\sum_b \frac{\hat\sigma_c\, r^{(c)}_b\, V'^{(c)}_b}
            {V_\bullet^{(c)} + \hat\sigma_c\, r^{(c)}_b}
= U_\bullet^{(c)},
\qquad
\hat\eqm^{(c)}_b = \frac{V'^{(c)}_b}{V_\bullet^{(c)} + \hat\sigma_c\, r^{(c)}_b},
\end{equation}
with $r^{(c)}_b = \sum_{a\neq b}\exch^{(c)}_{ab} W^{(c)}_a$.  Classes
do not couple, by the same class-decomposition~\eqref{eq:mix-loglik}
that justifies per-class M-steps in the standard mixture regime.

\paragraph{Aggregation of counts.}
The required sufficient statistics per class are
$(U_\bullet^{(c)}, V_\bullet^{(c)}, V'^{(c)}, r^{(c)})$.  All are
linear contractions of the underlying per-pair $(W, U, V)$ statistics
against the fixed quantities $(\exch^{(c)}, V_\bullet^{(c)} = $ count
of $V$ entries$)$, so under SVI--BW they accumulate via the same EMA
used for the standard M-step (no special handling).

\paragraph{Reduction to pure rescaling.}
Setting $\eqm = \eqm_{\mathrm{old}}$ in \eqref{eq:rescale-pi-foc-sigma}
recovers the closed form $\hat\sigma = U_\bullet / D$ of
\eqref{eq:rescale-sigma}.  The joint mode strictly dominates pure
rescaling whenever $\hat\eqm \neq \eqm_{\mathrm{old}}$ at the joint
optimum (i.e.\ whenever the Maraschino-fit warm $\eqm$ is not exactly
the M-step optimum given the current $W, U, V$); in practice the gain
is largest immediately after a warm start that fixed $\eqm$ away from
its joint optimum.

\paragraph{When to use.}
Use the joint $(\sigma, \eqm)$ rescaling when warm-starting from a
fit that already learnt per-class $(\sigma_c, \eqm^{(c)})$ but the
SVI--BW M-step should keep updating both with new evidence.  Use the
pure $\sigma$-only rescaling when the warm $\eqm^{(c)}$ should be
treated as a strong prior or when one wishes to attribute any
substitution-rate drift purely to the time scale.


%%%%%%%%%%%%%%%%%%%%%%%%%%%%%%%%%%%%%%%%%%%%%%%%%%%%%%%%%%%%%%%%%%%%%%%%
\subsection{Tied Equilibria across Class Blocks}
\label{sec:mstep-tied-pi}

A second restricted regime ties the equilibrium distribution
$\eqm^{(c)}$ across blocks of site classes while leaving
exchangeabilities $\exch^{(c)}$ free.  Partition the $C$ classes
$\{1,\dots,C\}$ into $B = C / N_{\mathrm{tied}}$ disjoint blocks
$\mathcal{B}_1,\dots,\mathcal{B}_B$ each of size $N_{\mathrm{tied}}$
(we require $N_{\mathrm{tied}} \mid C$).  Within each block~$b$ all
classes share a common equilibrium distribution
\begin{equation}\label{eq:tied-pi-tying}
\eqm^{(c)} \equiv \tilde\eqm^{(b)}
\quad \text{for all } c \in \mathcal{B}_b.
\end{equation}

\paragraph{Block-pooled complete-data log-likelihood.}
Substituting~\eqref{eq:tied-pi-tying} into~\eqref{eq:mix-loglik}, the
log-likelihood decomposes across blocks rather than across individual
classes:
\begin{equation}\label{eq:tied-pi-loglik}
\ell_2 = \sum_{b=1}^{B} \ell_2^{(b)},
\end{equation}
\begin{align}
\ell_2^{(b)} &=
   \sum_i \tilde V'^{(b)}_i \log \tilde\eqm^{(b)}_i
   \;+\; \sum_{c\in\mathcal{B}_b}\sum_{j>i}
        \bigl(U^{(c)}_{ij}+U^{(c)}_{ji}\bigr)\log\exch^{(c)}_{ij}
   \notag\\
&\quad - \sum_{c\in\mathcal{B}_b}\sum_{j>i}
   \exch^{(c)}_{ij}\,
   \bigl(W^{(c)}_i\,\tilde\eqm^{(b)}_j
        + W^{(c)}_j\,\tilde\eqm^{(b)}_i\bigr),
\label{eq:tied-pi-loglik-b}
\end{align}
with the pooled equilibrium-character count
\begin{equation}\label{eq:tied-pi-Vprime}
\tilde V'^{(b)}_i = \sum_{c\in\mathcal{B}_b} V'^{(c)}_i.
\end{equation}
The exchangeability parts~$\exch^{(c)}$ remain class-specific; only
the equilibrium term pools.

\paragraph{Block M-step.}
Within block~$b$, alternate between the per-class $\exch^{(c)}$ update
(holding $\tilde\eqm^{(b)}$ fixed) and the pooled $\tilde\eqm^{(b)}$
update (holding all $\exch^{(c)}, c\in\mathcal{B}_b$ fixed).
The $\exch$-step uses the standard mixture
formula~\eqref{eq:mix-S-update} with the shared $\tilde\eqm^{(b)}$:
\begin{equation}\label{eq:tied-pi-S-update}
\hat\exch^{(c)}_{ij}
= \frac{U^{(c)}_{ij}+U^{(c)}_{ji}}
       {W^{(c)}_i\,\tilde\eqm^{(b)}_j
       + W^{(c)}_j\,\tilde\eqm^{(b)}_i}
\quad (c\in\mathcal{B}_b).
\end{equation}
The $\tilde\eqm^{(b)}$-step is the same Lagrange-multiplier solve as
the unconstrained mixture M-step, but with pooled coefficients.
Differentiating~\eqref{eq:tied-pi-loglik-b}:
\begin{equation}\label{eq:tied-pi-pi-update}
\tilde\eqm^{(b)}_i
= \frac{\tilde V'^{(b)}_i}{\tilde c^{(b)}_i - \eta},
\qquad
\tilde c^{(b)}_i
= \sum_{c\in\mathcal{B}_b}\sum_{j\neq i}\exch^{(c)}_{ij}\,W^{(c)}_j,
\end{equation}
with $\eta$ found by 1D bisection from
$\sum_i \tilde V'^{(b)}_i / (\tilde c^{(b)}_i - \eta) = 1$, exactly as
in Section~\ref{sec:bw-tkf91}.  Each coordinate update is an exact
maximiser within its coordinate; the joint objective on
$(\exch^{(\cdot)}, \tilde\eqm^{(b)})$ is strictly concave on
$\mathbb{R}_+^{|\mathcal{B}_b| \cdot \binom{|\alphabet|}{2}}
   \times \Delta^{|\alphabet|-1}$ (the latter the simplex), so the
coordinate ascent converges to the unique block MLE.

\paragraph{Connection to standard models.}
With $N_{\mathrm{tied}} = 1$ (one class per block) the tying is
trivial and~\eqref{eq:tied-pi-pi-update} reduces to the per-class
Lagrange solve of Section~\ref{sec:mstep-mixture}.
With $N_{\mathrm{tied}} = C$ (one block containing all classes) all
classes share a single $\tilde\eqm$ — the discrete-rate-classes
parameterisation of~\cite{Yang1994a} when the $\exch^{(c)}$ are also
proportional to a shared shape.
For intermediate $N_{\mathrm{tied}}$ the tying captures shared
equilibrium chemistry across classes that nonetheless differ in their
exchangeability (and hence in their substitution \emph{pattern} and
\emph{rate}), without forcing every class to maintain its own
equilibrium estimate.

\paragraph{Pseudocounts.}
The LG-informed Dirichlet pseudocount on~$\tilde\eqm^{(b)}$ adds
$N_\pi\,\eqm^{\mathrm{LG}}_i$ to $\tilde V'^{(b)}_i$ once per block
(not once per class within the block); the per-class
Gamma$(a_S, b_S)$ pseudocount on~$\exch^{(c)}_{ij}$ enters the
class-specific update~\eqref{eq:tied-pi-S-update} unchanged.

\paragraph{Combination with rate rescaling.}
Tied-$\eqm$ and rate rescaling can also be \emph{combined} within a
single M-step: each class has its own scalar $\sigma_c$ but classes
within a block share the equilibrium $\tilde\eqm^{(b)}$.  The
coordinate-ascent factorisation does dispatch in that joint regime;
see Section~\ref{sec:mstep-tied-pi-rescaling}.  The frozen-$\eqm$
mode (which fixes $\eqm^{(c)}$ at its initialisation, with no M-step
update) remains an alternative restriction.  Alternation
\emph{across} M-steps (rate rescaling on odd iterations, tied-$\eqm$
on even iterations) is valid in addition because each step is an
unconditional ascent on the joint $\ell_2$.


%%%%%%%%%%%%%%%%%%%%%%%%%%%%%%%%%%%%%%%%%%%%%%%%%%%%%%%%%%%%%%%%%%%%%%%%
\subsection{Tied Equilibria with Per-Class Rate Rescaling}
\label{sec:mstep-tied-pi-rescaling}

The two restrictions of Sections~\ref{sec:mstep-rescaling} and
\ref{sec:mstep-tied-pi} can be applied jointly to the same mixture
M-step: each class~$c$ in block~$\mathcal{B}_b$ has its own
\emph{scalar} rate multiplier $\sigma_c$, with the exchangeability
shape $\exch^{(c)}$ frozen, while all classes within the block share
a common equilibrium $\tilde\eqm^{(b)}$.  This matches the
substitution degrees of freedom of a Maraschino-style fit with
\verb|--freeze-class-S-shape| under a block-tied
\verb|class_pi| constraint and is strictly more constrained than the
joint $(\sigma_c, \eqm^{(c)})$ M-step of
Section~\ref{sec:mstep-rescaling-pi}.

\paragraph{Block-pooled complete-data log-likelihood.}
With $\revsub^{(c)}_{ij} = \sigma_c\,\exch^{(c)}_{ij}\,\tilde\eqm^{(b)}_j$
for $c \in \mathcal{B}_b$ (with $\exch^{(c)}$ frozen), the per-block
log-likelihood is
\begin{equation}\label{eq:tpr-ll}
\ell_2^{(b)}(\{\sigma_c\}_{c\in\mathcal{B}_b}, \tilde\eqm^{(b)}) =
   \sum_b' \tilde V'^{(b)}_{b'} \log\tilde\eqm^{(b)}_{b'}
   + \sum_{c\in\mathcal{B}_b}\Bigl[ U^{(c)}_\bullet \log\sigma_c
       - \sigma_c\, \mathbf{r}^{(c)\top} \tilde\eqm^{(b)} \Bigr]
   + \mathrm{const},
\end{equation}
where the constant collects $\sum_c \sum_{j>i}(U^{(c)}_{ij}+U^{(c)}_{ji})
\log\exch^{(c)}_{ij}$ (frozen),
$\tilde V'^{(b)}_{b'} = \sum_{c\in\mathcal{B}_b} V'^{(c)}_{b'}$
is the block-pooled boundary count, and
$r^{(c)}_{b'} = \sum_{a\neq b'} \exch^{(c)}_{a,b'} W^{(c)}_a$
depends only on the frozen $(\exch^{(c)}, W^{(c)})$.

\paragraph{Stationary conditions.}
Lagrangian for $\sum_{b'}\tilde\eqm^{(b)}_{b'} = 1$ with multiplier
$\lambda^{(b)}$:
\begin{align}
\frac{\partial}{\partial \sigma_c}
   &= \frac{U^{(c)}_\bullet}{\sigma_c}
        - \mathbf{r}^{(c)\top}\tilde\eqm^{(b)} = 0
   \;\Longrightarrow\;
   \sigma_c = \frac{U^{(c)}_\bullet}
                   {\mathbf{r}^{(c)\top}\tilde\eqm^{(b)}},
   \label{eq:tpr-foc-sigma}\\[2pt]
\frac{\partial}{\partial \tilde\eqm^{(b)}_{b'}}
   &= \frac{\tilde V'^{(b)}_{b'}}{\tilde\eqm^{(b)}_{b'}}
        - \sum_{c\in\mathcal{B}_b} \sigma_c\,r^{(c)}_{b'}
        - \lambda^{(b)} = 0
   \;\Longrightarrow\;
   \tilde\eqm^{(b)}_{b'} = \frac{\tilde V'^{(b)}_{b'}}
       {\lambda^{(b)} + \sum_{c\in\mathcal{B}_b} \sigma_c r^{(c)}_{b'}}.
   \label{eq:tpr-foc-pi}
\end{align}

\paragraph{Block Lagrange multiplier is fixed.}
Multiplying~\eqref{eq:tpr-foc-pi} by $\tilde\eqm^{(b)}_{b'}$ and
summing over $b'$,
\[
\sum_{b'} \tilde V'^{(b)}_{b'}
= \lambda^{(b)} + \sum_{c\in\mathcal{B}_b} \sigma_c\,
     \mathbf{r}^{(c)\top}\tilde\eqm^{(b)}.
\]
Substituting~\eqref{eq:tpr-foc-sigma} into the right-hand side
collapses the $\sigma_c$-dependent term to $\sum_c U^{(c)}_\bullet$,
giving
\begin{equation}\label{eq:tpr-lambda}
\boxed{\;\lambda^{(b)} \;=\; \tilde V_\bullet^{(b)},\;}
\qquad
\tilde V_\bullet^{(b)} := \sum_{c\in\mathcal{B}_b} V_\bullet^{(c)}
= \sum_{c\in\mathcal{B}_b}\sum_i V^{(c)}_i,
\end{equation}
generalising the single-class result~\eqref{eq:rescale-pi-lambda}.
The multiplier is fully determined by the block-pooled boundary
count.

\paragraph{Closed-form coupled updates.}
Substituting~\eqref{eq:tpr-lambda} into~\eqref{eq:tpr-foc-pi},
\begin{equation}\label{eq:tpr-pi-update}
\hat{\tilde\eqm}^{(b)}_{b'} = \frac{\tilde V'^{(b)}_{b'}}
   {\tilde V_\bullet^{(b)}
    + \sum_{c\in\mathcal{B}_b} \sigma_c\, r^{(c)}_{b'}}.
\end{equation}
Equation~\eqref{eq:tpr-pi-update} together with
\eqref{eq:tpr-foc-sigma} forms a closed nonlinear system in the
$|\mathcal{B}_b|+A$ unknowns $(\{\sigma_c\}, \tilde\eqm^{(b)})$ per
block.  Two practical solvers, both monotone in $\ell_2^{(b)}$:

\begin{enumerate}
\item \emph{Coordinate ascent.}  Alternate
  $\sigma_c \leftarrow U^{(c)}_\bullet / (\mathbf{r}^{(c)\top}\tilde\eqm^{(b)})$
  for each $c \in \mathcal{B}_b$ and
  $\tilde\eqm^{(b)} \leftarrow \tilde V'^{(b)} \,/\,
  (\tilde V_\bullet^{(b)} + \sum_c \sigma_c r^{(c)})$ (componentwise),
  with a final renormalisation safety guard.  Each update is the
  exact closed-form maximiser in its own coordinate; convergence to
  the unique block-optimum follows in $\sim 5\text{--}20$ rounds.
\item \emph{Substitute and root-find.}  Eliminate $\sigma_c$ from
  \eqref{eq:tpr-pi-update} via~\eqref{eq:tpr-foc-sigma}; the result is
  a vector fixed-point equation in $\tilde\eqm^{(b)} \in \Delta^{A-1}$
  alone.  When $|\mathcal{B}_b|$ is small (typical: 2--5), the
  resulting low-dimensional system is amenable to a damped Newton
  iteration on the $A-1$ free $\tilde\eqm$ components; in practice
  coordinate ascent is faster and simpler.
\end{enumerate}
Both solvers reach the same fixed point because $\ell_2^{(b)}$ is
strictly concave in each block of variables given the others (the
$\log$ terms ensure positivity through the FOCs).

\paragraph{Mixture-wide M-step.}
The blocks decouple in the same way the classes do under
Section~\ref{sec:mstep-tied-pi}: the joint
$\ell_2 = \sum_b \ell_2^{(b)}$ separates over $b$, so each block can
be solved independently.

\paragraph{Aggregation of counts.}
Per block, the required sufficient statistics are the per-class
$(U^{(c)}_\bullet, V'^{(c)}, r^{(c)})$ for $c \in \mathcal{B}_b$ plus
the block totals $\tilde V'^{(b)}, \tilde V_\bullet^{(b)}$.  All are
linear in $(W, U, V)$ and accumulate via the same EMA used by the
standard M-step.

\paragraph{Reduction to special cases.}
Setting $|\mathcal{B}_b| = 1$ (each class its own block) recovers
the joint $(\sigma_c, \eqm^{(c)})$ M-step of
Section~\ref{sec:mstep-rescaling-pi}; setting all $\sigma_c = 1$
(no rate rescaling within block) recovers a degenerate version of
the tied-$\eqm$ M-step with frozen exchangeabilities; freezing
$\tilde\eqm^{(b)}$ at its current value recovers the rate-only
rescaling of Section~\ref{sec:mstep-rescaling}.  Strict dominance
holds in both directions: this regime improves on
Section~\ref{sec:mstep-rescaling} by adding the $\eqm$ degree of
freedom and improves on Section~\ref{sec:mstep-tied-pi} (with
$\exch^{(c)}$ frozen at LG) by adding per-class $\sigma_c$.

\paragraph{When to use.}
Use this regime when the warm-start fit (e.g.\ Maraschino with
\verb|--freeze-class-S-shape|) ties classes into blocks with shared
equilibrium chemistry but distinct rates, and the SVI--BW M-step
should preserve that tying while continuing to update both
$(\sigma_c, \tilde\eqm^{(b)})$ from new evidence.  The block tying
is a strong inductive bias when the number of classes per block is
small relative to $A$ and the per-class equilibrium counts
$V'^{(c)}$ would otherwise be too sparse to estimate
$\eqm^{(c)}$ reliably.

